% GENERATED FILE. Edit canonical lessons, then run npm run generate:books.
% canonical-source-hash: fnv1a64:66d115e53417dc33
% canonical-lessons: GU-C505-kapas, GU-C505-ru, GU-C505-retiyo, GU-C505-katvu, GU-C505-khadi

\chapter{Cotton and the Spinning Wheel}
\label{ch:gu-cotton-and-the-spinning-wheel}
\hlchaptermodality{505}

\begin{chapteropening}
\textbf{By the end of this chapter:} \emph{I can talk about cotton, spinning and handspun cloth.}

Complete the last lesson of chapter 505: I can talk about cotton, spinning and handspun cloth.
\end{chapteropening}

\section[\texorpdfstring{\gu{કપાસ} (kapās)}{કપાસ (kapās)}]{\gu{કપાસ} (kapās) — the cotton plant}

\label{lesson:GU-C505-kapas}



\begin{quote}
Before the new one: say the Gujarati for to knead clay, then the Gujarati for a clay lamp.
\end{quote}

\subsection*{You'll want to know: \gu{કપાસ}}
\textbf{\gu{કપાસ}} — \emph{kapās} — \textquotedblleft{}the cotton plant\textquotedblright{}.

\begin{quote}
\gu{અમારા} \gu{ગામનાં} \gu{ખેતરોમાં} \gu{ઘણો} \gu{કપાસ} \gu{ઊગે} \gu{છે}.
\end{quote}

\begin{quote}
\emph{amārā gāmnā̃ khetromā̃ ghaṇo kapās ūge che.}
\end{quote}

\begin{quote}
A lot of cotton grows in our village's fields.
\end{quote}

A masculine noun.

\begin{cousinweb}
\textbf{\gu{કપાસ}} (\emph{kapās}) is a tadbhava, inherited from Sanskrit \emph{karpāsa}, \textquotedblleft{}the cotton plant\textquotedblright{}, and CDIAL lists the Gujarati word (CDIAL 2877).
\end{cousinweb}

\begin{grammarlens}[title={What you've built}]
One more word for describing what you lived and felt.
\end{grammarlens}

\subsection*{Your turn}
\begin{itemize}
\raggedright
  \item \emph{Say it:} \emph{kapās}
  \item \emph{Say it:} \emph{kapās}, once more
  \item \emph{Say it:} \emph{amārā gāmnā̃ khetromā̃ ghaṇo kapās ūge che.} — A lot of cotton grows in our village's fields.
  \item \emph{Recall:} say the Gujarati for to knead clay, then the Gujarati for a clay lamp, then say \emph{kapās} again
\end{itemize}

\begin{tcolorbox}[breakable,colback=teal!4,colframe=teal!35!black,title={Before you move on}]
What does \gu{કપાસ} mean? (\textquotedblleft{}The cotton plant\textquotedblright{}.) Say it once more.
\end{tcolorbox}
\section[\texorpdfstring{\gu{રૂ} (rū)}{રૂ (rū)}]{\gu{રૂ} (rū) — raw cotton (for spinning)}

\label{lesson:GU-C505-ru}



\begin{quote}
Before the new one: say the Gujarati for a clay lamp, then the Gujarati for the cotton plant.
\end{quote}

\subsection*{You'll want to know: \gu{રૂ}}
\textbf{\gu{રૂ}} — \emph{rū} — \textquotedblleft{}raw cotton (for spinning)\textquotedblright{}.

\begin{quote}
\gu{મારી} \gu{દાદી} \gu{રૂમાંથી} \gu{દોરો} \gu{બનાવતી} \gu{હતી}.
\end{quote}

\begin{quote}
\emph{mārī dādī rūmā̃thī doro banāvtī hatī.}
\end{quote}

\begin{quote}
My grandmother used to make thread from cotton.
\end{quote}

A neuter noun.

\begin{cousinweb}
\textbf{\gu{રૂ}} (\emph{rū}) comes from \emph{*rū-a-}, \textquotedblleft{}cotton\textquotedblright{}, a form CDIAL reconstructs, and CDIAL lists the Gujarati word as carded cotton (CDIAL 10798).
\end{cousinweb}

\begin{grammarlens}[title={What you've built}]
One more word for describing what you lived and felt.
\end{grammarlens}

\subsection*{Your turn}
\begin{itemize}
\raggedright
  \item \emph{Say it:} \emph{rū}
  \item \emph{Say it:} \emph{rū}, once more
  \item \emph{Say it:} \emph{mārī dādī rūmā̃thī doro banāvtī hatī.} — My grandmother used to make thread from cotton.
  \item \emph{Recall:} say the Gujarati for a clay lamp, then the Gujarati for the cotton plant, then say \emph{rū} again
\end{itemize}

\begin{tcolorbox}[breakable,colback=teal!4,colframe=teal!35!black,title={Before you move on}]
What does \gu{રૂ} mean? (\textquotedblleft{}Raw cotton (for spinning)\textquotedblright{}.) Say it once more.
\end{tcolorbox}
\section[\texorpdfstring{\gu{રેંટિયો} (rẽṭiyo)}{રેંટિયો (rẽṭiyo)}]{\gu{રેંટિયો} (rẽṭiyo) — a spinning wheel}

\label{lesson:GU-C505-retiyo}



\begin{quote}
Before the new one: say the Gujarati for the cotton plant, then the Gujarati for raw cotton.
\end{quote}

\subsection*{You'll want to know: \gu{રેંટિયો}}
\textbf{\gu{રેંટિયો}} — \emph{rẽṭiyo} — \textquotedblleft{}a spinning wheel\textquotedblright{}.

\begin{quote}
\gu{દાદીનો} \gu{રેંટિયો} \gu{હજી} \gu{અમારા} \gu{ઘરમાં} \gu{છે}.
\end{quote}

\begin{quote}
\emph{dādīno rẽṭiyo hajī amārā gharmā̃ che.}
\end{quote}

\begin{quote}
Grandmother's spinning wheel is still in our house.
\end{quote}

A masculine noun.

\begin{cousinweb}
\textbf{\gu{રેંટિયો}} (\emph{rẽṭiyo}) is filed by CDIAL under Sanskrit \emph{araghaṭṭa}, \textquotedblleft{}a wheel for raising water\textquotedblright{}, and CDIAL lists the Gujarati word as a spinning wheel (CDIAL 596): one turning wheel gave its name to another.
\end{cousinweb}

\begin{grammarlens}[title={What you've built}]
One more word for describing what you lived and felt.
\end{grammarlens}

\subsection*{Your turn}
\begin{itemize}
\raggedright
  \item \emph{Say it:} \emph{rẽṭiyo}
  \item \emph{Say it:} \emph{rẽṭiyo}, once more
  \item \emph{Say it:} \emph{dādīno rẽṭiyo hajī amārā gharmā̃ che.} — Grandmother's spinning wheel is still in our house.
  \item \emph{Recall:} say the Gujarati for the cotton plant, then the Gujarati for raw cotton, then say \emph{rẽṭiyo} again
\end{itemize}

\begin{tcolorbox}[breakable,colback=teal!4,colframe=teal!35!black,title={Before you move on}]
What does \gu{રેંટિયો} mean? (\textquotedblleft{}A spinning wheel\textquotedblright{}.) Say it once more.
\end{tcolorbox}
\section[\texorpdfstring{\gu{કાંતવું} (kā̃tvũ)}{કાંતવું (kā̃tvũ)}]{\gu{કાંતવું} (kā̃tvũ) — to spin (thread)}

\label{lesson:GU-C505-katvu}



\begin{quote}
Before the new one: say the Gujarati for raw cotton, then the Gujarati for a spinning wheel.
\end{quote}

\subsection*{You'll want to know: \gu{કાંતવું}}
\textbf{\gu{કાંતવું}} — \emph{kā̃tvũ} — \textquotedblleft{}to spin (thread)\textquotedblright{}.

\begin{quote}
\gu{દાદી} \gu{રેંટિયા} \gu{પર} \gu{રૂ} \gu{કાંતતી} \gu{હતી}.
\end{quote}

\begin{quote}
\emph{dādī rẽṭiyā par rū kā̃ttī hatī.}
\end{quote}

\begin{quote}
Grandmother used to spin cotton on the spinning wheel.
\end{quote}

Past, with \gu{તેણે} (he or she …): \gu{તેણે} \gu{દોરો} \gu{કાંત્યો}, \emph{teṇe doro kā̃tyo}.

\begin{cousinweb}
\textbf{\gu{કાંતવું}} (\emph{kā̃tvũ}) is inherited from Sanskrit \emph{kṛntati}, \textquotedblleft{}spins\textquotedblright{}, and CDIAL lists the Gujarati verb (CDIAL 3433).
\end{cousinweb}

\begin{grammarlens}[title={What you've built}]
One more word for describing what you lived and felt.
\end{grammarlens}

\subsection*{Your turn}
\begin{itemize}
\raggedright
  \item \emph{Say it:} \emph{kā̃tvũ}
  \item \emph{Say it:} \emph{kā̃tvũ}, once more
  \item \emph{Say it:} \emph{dādī rẽṭiyā par rū kā̃ttī hatī.} — Grandmother used to spin cotton on the spinning wheel.
  \item \emph{Recall:} say the Gujarati for raw cotton, then the Gujarati for a spinning wheel, then say \emph{kā̃tvũ} again
\end{itemize}

\begin{tcolorbox}[breakable,colback=teal!4,colframe=teal!35!black,title={Before you move on}]
What does \gu{કાંતવું} mean? (\textquotedblleft{}To spin (thread)\textquotedblright{}.) Say it once more.
\end{tcolorbox}
\section[\texorpdfstring{\gu{ખાદી} (khādī)}{ખાદી (khādī)}]{\gu{ખાદી} (khādī) — khadi (handspun cloth)}

\label{lesson:GU-C505-khadi}



\begin{quote}
Before the new one: say the Gujarati for a spinning wheel, then the Gujarati for to spin (thread).
\end{quote}

\subsection*{You'll want to know: \gu{ખાદી}}
\textbf{\gu{ખાદી}} — \emph{khādī} — \textquotedblleft{}khadi (handspun cloth)\textquotedblright{}.

\begin{quote}
\gu{મીરાએ} \gu{ખાદીનો} \gu{કુરતો} \gu{ખરીદ્યો}.
\end{quote}

\begin{quote}
\emph{Mīrāe khādīno kurto kharīdyo.}
\end{quote}

\begin{quote}
Mira bought a khadi kurta.
\end{quote}

A feminine noun.

\begin{cousinweb}
\textbf{\gu{ખાદી}} (\emph{khādī}) comes from \emph{*khadda}, \textquotedblleft{}coarse cloth\textquotedblright{}, a form CDIAL reconstructs, and CDIAL lists the Gujarati word as thick coarse cotton cloth (CDIAL 3808).
\end{cousinweb}

\begin{grammarlens}[title={What you've built}]
One more word for describing what you lived and felt.
\end{grammarlens}

\subsection*{Your turn}
\begin{itemize}
\raggedright
  \item \emph{Say it:} \emph{khādī}
  \item \emph{Say it:} \emph{khādī}, once more
  \item \emph{Say it:} \emph{Mīrāe khādīno kurto kharīdyo.} — Mira bought a khadi kurta.
  \item \emph{Recall:} say the Gujarati for a spinning wheel, then the Gujarati for to spin (thread), then say \emph{khādī} again
\end{itemize}

\begin{tcolorbox}[breakable,colback=teal!4,colframe=teal!35!black,title={Before you move on}]
What does \gu{ખાદી} mean? (\textquotedblleft{}Khadi (handspun cloth)\textquotedblright{}.) Say it once more.
\end{tcolorbox}
